% Antecedentes: estudios previos que sustentan el proyecto.
% Módulo independiente; todavía no se carga en ningún documento. Para usarlo,
% agregar % Antecedentes: estudios previos que sustentan el proyecto.
% Módulo independiente; todavía no se carga en ningún documento. Para usarlo,
% agregar % Antecedentes: estudios previos que sustentan el proyecto.
% Módulo independiente; todavía no se carga en ningún documento. Para usarlo,
% agregar % Antecedentes: estudios previos que sustentan el proyecto.
% Módulo independiente; todavía no se carga en ningún documento. Para usarlo,
% agregar \input{secciones/11-antecedentes} donde corresponda en trabajo.tex.
\section{Antecedentes}
El tiempo que esperan los estudiantes para ser atendidos puede analizarse
con la teoría de colas, a partir de la tasa de llegadas y del tiempo de
atención. Este enfoque ya se ha propuesto para reducir los tiempos de espera
de estudiantes en el área de control escolar de otro plantel del Tecnológico
Nacional de México \parencite{talavera2025espera}, lo que respalda su uso
para estudiar las filas que se forman en el acceso del \tescha{}.

En cuanto a la tecnología, el uso de códigos QR para el control de acceso
físico se ha planteado en la literatura \parencite{kao2011physical}. En una
universidad del Reino Unido, un código QR que se renovaba cada 20 segundos
impidió que los estudiantes reenviaran capturas de pantalla para registrar
su asistencia \parencite{nwabuwe2023dynamicqr}. Si además el código se firma
con una contraseña de un solo uso basada en HMAC, resiste la interceptación,
el reenvío y la modificación del mensaje \parencite{subpratatsavee2015hmac}.
 donde corresponda en trabajo.tex.
\section{Antecedentes}
El tiempo que esperan los estudiantes para ser atendidos puede analizarse
con la teoría de colas, a partir de la tasa de llegadas y del tiempo de
atención. Este enfoque ya se ha propuesto para reducir los tiempos de espera
de estudiantes en el área de control escolar de otro plantel del Tecnológico
Nacional de México \parencite{talavera2025espera}, lo que respalda su uso
para estudiar las filas que se forman en el acceso del \tescha{}.

En cuanto a la tecnología, el uso de códigos QR para el control de acceso
físico se ha planteado en la literatura \parencite{kao2011physical}. En una
universidad del Reino Unido, un código QR que se renovaba cada 20 segundos
impidió que los estudiantes reenviaran capturas de pantalla para registrar
su asistencia \parencite{nwabuwe2023dynamicqr}. Si además el código se firma
con una contraseña de un solo uso basada en HMAC, resiste la interceptación,
el reenvío y la modificación del mensaje \parencite{subpratatsavee2015hmac}.
 donde corresponda en trabajo.tex.
\section{Antecedentes}
El tiempo que esperan los estudiantes para ser atendidos puede analizarse
con la teoría de colas, a partir de la tasa de llegadas y del tiempo de
atención. Este enfoque ya se ha propuesto para reducir los tiempos de espera
de estudiantes en el área de control escolar de otro plantel del Tecnológico
Nacional de México \parencite{talavera2025espera}, lo que respalda su uso
para estudiar las filas que se forman en el acceso del \tescha{}.

En cuanto a la tecnología, el uso de códigos QR para el control de acceso
físico se ha planteado en la literatura \parencite{kao2011physical}. En una
universidad del Reino Unido, un código QR que se renovaba cada 20 segundos
impidió que los estudiantes reenviaran capturas de pantalla para registrar
su asistencia \parencite{nwabuwe2023dynamicqr}. Si además el código se firma
con una contraseña de un solo uso basada en HMAC, resiste la interceptación,
el reenvío y la modificación del mensaje \parencite{subpratatsavee2015hmac}.
 donde corresponda en trabajo.tex.
\section{Antecedentes}
El tiempo que esperan los estudiantes para ser atendidos puede analizarse
con la teoría de colas, a partir de la tasa de llegadas y del tiempo de
atención. Este enfoque ya se ha propuesto para reducir los tiempos de espera
de estudiantes en el área de control escolar de otro plantel del Tecnológico
Nacional de México \parencite{talavera2025espera}, lo que respalda su uso
para estudiar las filas que se forman en el acceso del \tescha{}.

En cuanto a la tecnología, el uso de códigos QR para el control de acceso
físico se ha planteado en la literatura \parencite{kao2011physical}. En una
universidad del Reino Unido, un código QR que se renovaba cada 20 segundos
impidió que los estudiantes reenviaran capturas de pantalla para registrar
su asistencia \parencite{nwabuwe2023dynamicqr}. Si además el código se firma
con una contraseña de un solo uso basada en HMAC, resiste la interceptación,
el reenvío y la modificación del mensaje \parencite{subpratatsavee2015hmac}.
